\chapter{Stereo Image and Panning}
\label{chap:stereo-image}

\classquote{有无相生，难易相成，长短相形，高下相倾，音声相和，前后相随。}{《老子·第二章》 / Laozi, Chapter 2}{Translation: Being and non-being generate each other; long and short define each other; high and low lean on each other; sound and voice harmonize; front and back follow. Note: stereo is relational — left only means something against right.}

Pan answers ``where is this sound placed left or right'', stereo width answers ``how far apart are its two sides'', and mono compatibility answers ``how much remains after the left and right channels are summed''. These three are often collapsed into one ``image'' knob, but in Yeban they live on different layers: track pan at the source, width in the channel strip's M/S stage, low-frequency narrowing in the strip's mono-maker, and mono checking in the control-room monitor path. This chapter follows those four layers and marks the controls that are not provided.

\section{Three distinct layers}

Yeban does not have one global image control. It has three independent parameter families:

\begin{table}[htbp]
\centering
\small
\begin{tabular}{@{}p{0.17\linewidth}p{0.31\linewidth}p{0.40\linewidth}@{}}
\hline
Layer & Parameter / entry & Scope and boundary \\
\hline
Track & \texttt{TrackV3.pan}, each mixer strip's pan cell & $-1$--$+1$; every track, bus and master has its own; can be automated as \texttt{TrackPan} \\
Channel strip & \texttt{strip\_width}, \texttt{mono\_below\_hz}, \texttt{ms\_mid\_*} & Exists only when the track has a channel strip; not a width knob beside the fader \\
SFZ region & \texttt{pan}, \texttt{width}, \texttt{position} & Per sampler region; only in the SFZ file, no UI editor \\
\hline
\end{tabular}
\caption{The three layers of image parameters in Yeban}
\label{tab:stereo-layers}
\end{table}

Track pan moves the whole track left or right; channel-strip width spreads or collapses its left/right components; an SFZ region's pan places a particular slice of the sampled instrument. They can be combined, but one cannot do the job of another.

\begin{center}
\begin{tikzpicture}[
  node distance=7mm and 8mm,
  box/.style={draw, rounded corners=2pt, minimum width=22mm, minimum height=8mm, align=center, font=\scriptsize},
  branch/.style={draw, dashed, rounded corners=2pt, minimum width=24mm, minimum height=7mm, align=center, font=\scriptsize},
  arrow/.style={-{Latex[length=2mm]}, thick}
]
  \node[box] (src) {source: audio clip\\or SFZ region};
  \node[box, right=of src] (tpan) {track pan\\$\mathrm{pan}=-1..+1$};
  \node[box, right=of tpan] (strip) {channel strip\\width/mid/mono};
  \node[box, right=of strip] (bus) {bus and master\\their own pan};
  \node[box, right=of bus] (mon) {control room\\MONO/SIDE/DIM};
  \node[box, below=of mon] (out) {audio output\\export bypasses monitor};
  \node[branch, below=of strip] (sfz) {SFZ region image\\pan→width→position};
  \draw[arrow] (src) -- (tpan);
  \draw[arrow] (tpan) -- (strip);
  \draw[arrow] (strip) -- (bus);
  \draw[arrow] (bus) -- (mon);
  \draw[arrow] (mon) -- (out);
  \draw[arrow, dashed] (sfz) -- (tpan);
\end{tikzpicture}
\end{center}

\screenshot{mixer-pan-strip}{The mixer strip has a track selected; the pan cell shows C or L50/R50 and the fader and pan values are visible.}

\section{Track pan and the four pan laws}

Track pan is bounded to $-1$--$+1$: $-1$ is hard left, $0$ centre, $+1$ hard right. The UI converts the thousandths integer into readings such as \texttt{C}, \texttt{L50} and \texttt{R50}: centre shows \texttt{C}, negative values show \texttt{L} plus an integer percentage, positive values \texttt{R} plus a percentage. The mixer strip's pan cell, the inspector and the arrangement share one value.

Pan is not a simple left/right gain pair. It follows the project-level \texttt{audio\_config.pan\_law}. Yeban has four laws:

\begin{table}[htbp]
\centering
\small
\begin{tabular}{@{}p{0.30\linewidth}p{0.22\linewidth}p{0.20\linewidth}p{0.16\linewidth}@{}}
\hline
Serialized name & Curve & Curve centre (before the identity rule) & Centre level \\
\hline
\texttt{Linear} & $L=1-p$, $R=p$, $p=(pan+1)/2$ & $(0.5,0.5)$ & $-6.02$ dB \\
\texttt{ConstantPower\allowbreak Minus3dB} & $L=\cos\theta$, $R=\sin\theta$ & $(\sqrt2/2,\sqrt2/2)$ & $-3.01$ dB \\
\texttt{ConstantPower\allowbreak Minus4\_5dB} & $L=\cos^{\alpha}\theta$, $R=\sin^{\alpha}\theta$ & $(0.5957,0.5957)$ & $-4.50$ dB \\
\texttt{ConstantPower\allowbreak Minus6dB} & $L=\cos^2\theta$, $R=\sin^2\theta$ & $(0.5,0.5)$ & $-6.02$ dB \\
\hline
\end{tabular}
\caption{The four pan laws; $\theta=(pan+1)\pi/4$, $\alpha=0.45/\log_{10}2\approx1.4948677$}
\label{tab:pan-laws}
\end{table}

The default is \texttt{ConstantPower\allowbreak Minus3dB}: the sum of left and right energy is exactly one, so sweeping pan does not collapse the total energy. The linear law keeps the amplitude sum at one and loses 6 dB in the centre; $-4.5$ dB and $-6$ dB are the other two constant-power curves. The project field accepts the four serialized names above, but the current UI has no pan-law selector; change \texttt{audio\_config.pan\_law} directly in the project JSON, and note that MCP currently has no writer for it.

\subsection{Centre identity and the 3 dB step}

At dead centre Yeban does not use the curve centre. It uses identity:

\begin{center}
\texttt{centered\_pan\_gains(pan, law)}:\quad
$\texttt{pan}=0$ or non-finite $\Rightarrow (1.0,1.0)$; otherwise \texttt{pan\_gains}.
\end{center}

So a track reading \texttt{C} multiplies both sides by one, and monitoring and export agree bit-for-bit. As soon as pan leaves centre it falls onto the selected curve. For the default $-3$ dB law, $\mathrm{pan}=\pm\epsilon$ gives about $0.7071$ on each side, $3.01$ dB below centre; the step is $6.02$ dB for the linear law, about $4.50$ dB for $-4.5$ dB, and $6.02$ dB for $-6$ dB. This is Yeban's existing contract, not a typo: centre must align with the offline ``pan untouched'' case, and the price is an audible step between exact centre and a tiny non-zero pan.

The UI reading rounds to thousandths, so a very small non-zero pan may display \texttt{C} while processing already follows the curve. Do not infer ``identity'' from a displayed \texttt{C}; to reset exactly, write the value back to $0$.

\subsection{Tracks, buses and automation}

Tracks, auxiliary buses and the master each carry a \texttt{pan}, and all use the same \texttt{centered\_pan\_gains}. Master pan is not an output balance or a master gain; it moves the whole mix left or right. Do not use it to fix the position of one instrument.

Pan can be written as a \texttt{TrackPan} automation lane. When automation exists, it replaces the track's static \texttt{pan}, smoothed per sample with a 5 ms time constant; the static value is used only when no automation exists. See the automation chapter for drawing and reading; this chapter does not repeat it.

\section{Width: the channel strip's M/S scaling}

Yeban has no per-track width fader and no track parameter named \texttt{width}. Real width lives in the channel-strip device as \texttt{strip\_width} (alias \texttt{channel\_width}), range $0.0$--$2.0$, default $1.0$:

\begin{table}[htbp]
\centering
\small
\begin{tabular}{@{}llll@{}}
\hline
Parameter & Meaning & Range & Default \\
\hline
\texttt{strip\_width} & M/S width scaling; $1$ neutral, $0$ mono, $2$ double & 0.0--2.0 & 1.0 \\
\texttt{mono\_below\_hz} & mono-maker corner; $\le0$ bypass & 0--1000 Hz & 0 \\
\texttt{ms\_mid\_gain\_db} & peak gain on mid only & $-60$--$24$ dB & 0 \\
\texttt{ms\_mid\_freq\_hz} & M/S mid peak centre & 100--12000 Hz & 2000 \\
\hline
\end{tabular}
\caption{Channel-strip width and M/S parameters}
\label{tab:strip-width-params}
\end{table}

The width stage computes:

\begin{center}
\texttt{mid = 0.5(L+R)};\quad \texttt{side = 0.5(L-R) * strip\_width};\quad
\texttt{(L,R) = (mid+side, mid-side)}.
\end{center}

At $1.0$ the stage is skipped bit-for-bit: no M/S round trip, because in $f32$ $(m+s)+(m-s)$ is not guaranteed to equal $2m$. At $0.0$ side is multiplied by zero, both channels keep only mid, and the whole track becomes mono. At $2.0$ side is doubled, the left/right difference is exaggerated and the image feels wider, while the centre and low-frequency phase relationship is unchanged. Width never touches mid, so the centre of a vocal or kick does not drift sideways.

The width stage sits inside the channel strip: input gain → high-pass → EQ → dynamic EQ → filter → compressor → \textbf{width} → M/S mid → mono-maker → output gain. It exists only when the track has a channel strip; without one there is no \texttt{strip\_width}. The UI has no width knob; the device card shows recognised parameters and a curve only. To change width, use MCP \texttt{setParam} with \texttt{lane} set to \texttt{DeviceParam} and address \texttt{strip\_width} by name, or edit the device parameters in the project. The current version does \textbf{not provide} a track-level width knob, a dedicated automation lane named ``width'', or a width meter; automate width through the \texttt{DeviceParam} lane.

\screenshot{device-rack-strip-width}{F9 device rack with a channel-strip track selected; the card shows recognized parameter readouts and curves, with no width knob.}

\section{Low-frequency narrowing: mono-maker}

\texttt{mono\_below\_hz} is the channel strip's stage for narrowing low-frequency side energy, range $0$--$1000$ Hz, default $0$ (bypassed). Its contract is:

\begin{center}
\texttt{mid} unchanged;\quad \texttt{side} through a second-order Butterworth high-pass;\quad reconstructed to left/right.
\end{center}

It is not a brick-wall split: with a $120$ Hz corner, the side at $60$ Hz is about $-12$ dB, at $30$ Hz about $-24$ dB, and above $240$ Hz essentially untouched. It pins the low end to the centre; it does not make everything below $120$ Hz absolutely mono. For full-band mono, set \texttt{strip\_width} to $0$, not mono-maker.

Mono-maker runs after width: width scales the whole side (including low frequencies), then mono-maker removes the low-frequency side. Reversing the order leaves a different low-frequency residue. It also exists only inside a channel strip.

\section{M/S mid: centre only, sides untouched}

\texttt{ms\_mid\_gain\_db} is a peak EQ applied only to mid, range $-60$--$24$ dB, default $0$ (bypassed); \texttt{ms\_mid\_freq\_hz} ranges $100$--$12000$ Hz, default $2000$ Hz, with a fixed Q of $0.9$. It is not the ordinary EQ mid: \texttt{eq.mid\_*} is full-band, so boosting 3 kHz also boosts the cymbals and reverb tail in side. M/S mid boosts only the centre, bringing a lead vocal or melody forward without brightening the sides. At $0$ dB the stage is bypassed and performs no M/S round trip at all.

It is neither width nor a mono-compatibility tool. Boosting mid does not replace mono-maker; it makes the centre louder, not the low-frequency side disappear.

\section{The SFZ region image family}

An SFZ region can declare three image parameters, all in percent:

\begin{table}[htbp]
\centering
\small
\begin{tabular}{@{}llll@{}}
\hline
SFZ parameter & Meaning & Range & Default \\
\hline
\texttt{pan} & Region pan; placement for mono samples, balance for stereo & $-100$--$100$ & 0 \\
\texttt{width} & Stereo width; $100$ neutral, $0$ mono, negative swaps L/R & $-100$--$100$ & 100 \\
\texttt{position} & Moves the whole image after width & $-100$--$100$ & 0 \\
\hline
\end{tabular}
\caption{SFZ region image parameters}
\label{tab:sfz-pan-family}
\end{table}

The order is \texttt{pan} → \texttt{width} → \texttt{position}. \texttt{width} and \texttt{position} apply only to stereo samples; a mono sample's voice does not enter the split-channel path and the values are not applied. The default combination (\texttt{width=100}, \texttt{position=0}) short-circuits bit-for-bit and performs no floating-point work. All three are clamped to $-100$--$100$ at load time; out-of-range values raise a \texttt{ClampedStaticValue} warning.

Yeban can read mono and stereo sample files, but the UI has no SFZ region image editor; these values exist only in the SFZ file. The current version does \textbf{not provide} a graphical editor for region \texttt{pan}/\texttt{width}/\texttt{position}.

\section{Mono-compatibility checking}

Mono compatibility is not a plug-in; it is a set of switches in the control-room monitor path. Press F4 to show the control room and loudness meters. Its right side has \texttt{DIM}, \texttt{MONO} and \texttt{SIDE}; the lower-dock master strip also has a \texttt{MONO}, which is the same boolean as the control room, not a second state.

\screenshot{control-room-mono-side}{F4 control room; DIM/MONO/SIDE buttons are visible, with MONO or SIDE lit.}

\begin{table}[htbp]
\centering
\small
\begin{tabular}{@{}p{0.14\linewidth}p{0.30\linewidth}p{0.26\linewidth}p{0.22\linewidth}@{}}
\hline
Switch & Processing & Use & Boundary \\
\hline
\texttt{MONO} & both channels become $(L+R)/2$ & check cancellation and low-end thinning & monitor only, not export \\
\texttt{SIDE} & both channels become $L-R$ & hear the difference signal of reverb/wide sources & monitor only, not export \\
\texttt{DIM} & both channels multiply by $0.1$ ($-20$ dB) & temporarily lower monitoring level & monitor only, not export \\
\hline
\end{tabular}
\caption{Control-room monitor switches}
\label{tab:monitor-controls}
\end{table}

The order of the three switches is a contract: \textbf{SIDE first, then MONO, then DIM}. Therefore \texttt{SIDE} and \texttt{MONO} together still produce the difference signal, not silence; doing MONO first and SIDE second would subtract the two copies of $(L+R)/2$ and yield zero, which Yeban does not do. With all three off the stage is skipped entirely and touches no sample.

A checking routine:

\begin{enumerate}
  \item Monitor normally and note where the vocal, kick and bass sit in the centre.
  \item Press \texttt{MONO}; listen for a thinner vocal, a kick that loses low end, or reverb that suddenly grows. An instrument that almost disappears in mono is relying on left/right difference.
  \item Press \texttt{SIDE}; confirm the difference signal does not contain the dry centre you want to keep. \texttt{SIDE}+\texttt{MONO} together is still the difference signal, not silence.
  \item Turn \texttt{MONO}/\texttt{SIDE} off and return to normal monitoring; use \texttt{DIM} if you need a lower monitoring level.
\end{enumerate}

The current version does \textbf{not provide} automatic mono-compatibility analysis, a left/right correlation meter, a phase-cancellation warning, or a before/after mono-fold difference meter. \texttt{MONO} is a manual monitoring switch, not an analyser; it will not tell you ``this track lost 3 dB in mono''. The export path bypasses these switches, so do not judge export level while \texttt{MONO} is on.

\section{Common mistakes}

\begin{table}[htbp]
\centering
\small
\begin{tabular}{@{}p{0.38\linewidth}p{0.52\linewidth}@{}}
\hline
Mistake & Actual behaviour \\
\hline
``Pan is width'' & Pan moves the whole track; width scales only side, leaving mid alone \\
``Centre means 0 dB on the curve'' & Centre is identity at 0 dB; leaving centre drops onto the curve, a $3.01$ dB step for the default law \\
``Width 0 is the same as mono-maker at 120 Hz'' & Width 0 is full-band mono; mono-maker only rolls off low-frequency side at second order \\
``If an SFZ file says width, it must be audible'' & It applies only to stereo samples; mono samples ignore it \\
``MONO also makes the export mono'' & Monitor switches bypass export \\
``Channel-strip width can be dragged like a fader'' & No width knob exists; use device parameters/MCP or the project file \\
\hline
\end{tabular}
\caption{Common mistakes about pan and stereo}
\label{tab:stereo-mistakes}
\end{table}

\section{Exercise}


\begin{exercise}
\begin{enumerate}
  \item Move a track's pan from \texttt{C} to \texttt{L50}, then to \texttt{R50}; return to the full mix after each move and listen to the center and sides.
  \item Write the pan back to a true $0$ and confirm the display reads \texttt{C}. Do not infer identity processing from a ``C'' display.
  \item Check the track with control-room MONO/SIDE; record what is lost in mono, and do not use the monitoring level as the export level.
\end{enumerate}
\end{exercise}

\section{Help-topic mapping}

\begin{table}[htbp]
\centering
\small
\begin{tabular}{@{}p{0.34\linewidth}p{0.52\linewidth}@{}}
\hline
Section & In-app help anchor \\
\hline
Three layers & \texttt{stereo-layers} \\
Track pan and four laws & \texttt{track-pan-laws} \\
Centre identity and the step & \texttt{pan-center-identity} \\
Width and M/S parameters & \texttt{channel-strip-width} \\
mono-maker & \texttt{mono-maker} \\
SFZ region image family & \texttt{sfz-pan-family} \\
Mono-compatibility checking & \texttt{mono-compatibility} \\
Common mistakes & \texttt{stereo-mistakes} \\
\hline
\end{tabular}
\caption{Chapter 7 help-topic mapping}
\label{tab:stereo-help-map}
\end{table}

\section*{Further Reading}
\begin{itemize}[leftmargin=1.5em]
\item \url{https://www.izotope.com/en/learn.html} — iZotope Learn: study diagrams of panning, width, and mono compatibility, and fill in the left–right relationship.
\item \url{https://www.soundonsound.com/techniques} — Sound On Sound techniques: compare panning decisions in real mixes, and learn why wider is not automatically better.
\end{itemize}
